\documentclass{article}
\usepackage[submission]{colm2026_conference}

\usepackage{microtype}
\usepackage{hyperref}
\usepackage{url}
\usepackage{booktabs}
\usepackage{tabularx}
\usepackage{xcolor}
\usepackage{amsmath,amssymb}
\usepackage{lineno}

\definecolor{darkblue}{rgb}{0,0,0.5}
\hypersetup{
  colorlinks=true,
  linkcolor=darkblue,
  citecolor=darkblue,
  urlcolor=darkblue
}

\newcommand{\sys}{scientific AI}
\newcommand{\agents}{autonomous research agents}
\newcommand{\position}{accountable acceleration}
\newcommand{\safetycase}{scientific-agent safety case}

\title{Safety Cases for AI Scientists: \\
Accountable Acceleration in Language-Model-Driven Scientific Discovery}

\author{Anonymous Authors \\
Paper draft prepared for the LM4Sci 2026 Workshop}

\begin{document}

\ifcolmsubmission
\linenumbers
\fi

\maketitle

\begin{abstract}
Language models are moving from scientific assistants to research operators: they search literature, generate hypotheses, write code, run experiments, analyze data, and draft papers. The central governance question is therefore no longer whether AI can help produce science, but under what institutional conditions AI-produced science should be trusted, credited, and deployed. We argue that autonomous scientific AI creates an \emph{institutional agency gap}: current norms assign responsibility to human researchers, while emerging systems increasingly shape the selection of hypotheses, the execution of experiments, and the publication of claims. Generic model alignment, hallucination benchmarks, and vague human-in-the-loop review are insufficient for this gap because scientific agents act inside socio-technical systems that allocate resources, construct evidence, and affect public trust. We propose \position{} as a design goal for languague-model for scientific discovery (LM4Sci): accelerate discovery only when systems preserve epistemic integrity, safety, and accountability. Our position contributes (i) a ladder of scientific agency, (ii) a taxonomy of deployment risks specific to scientific AI, and (iii) a practical \safetycase{} framework built around capability boundaries, tool containment, provenance, human responsibility, and societal impact. We conclude with a research agenda for evaluating scientific agents as accountable research infrastructure rather than merely capable models.
\end{abstract}

\section{Introduction}

The language-model-for-science community is approaching a transition point. Early uses of language models in science emphasized text: summarizing papers, drafting code, extracting entities, or answering domain questions. The current frontier increasingly emphasizes \emph{agency}: systems that connect models to literature databases, code interpreters, simulation engines, robotic laboratories, theorem provers, molecule design tools, and manuscript generators. Language models can assist with idea generation, experiment execution, automated ML engineering, literature synthesis, data analysis, and even end-to-end research pipelines \citep{lm4sci2026cfp}. Recent systems illustrate the direction. Chemistry agents combine LLM planning with external chemical tools \citep{bran2023chemcrow}; autonomous chemistry systems have demonstrated planning and execution of experimental workflows \citep{boiko2023emergent}; systems such as \textsc{AI Scientist} automate the loop from ideas to experiments to papers \citep{lu2024aiscientist,yamada2025aiscientistv2}; biological and materials models such as AlphaFold and GNoME have shown how learned systems can reshape entire scientific subfields \citep{jumper2021alphafold,abramson2024alphafold3,merchant2023materials}. These examples do not imply that machines have replaced scientists. They do imply that the unit of analysis is changing from a model answering a prompt to an agent participating in a research institution.

This paper takes a position: \emph{scientific AI systems should be governed as accountable research infrastructure, not merely evaluated as language models}. The distinction matters. A language model can be wrong in a chat window; a scientific agent can select a bad hypothesis, run a misleading experiment, contaminate a benchmark, overuse scarce compute, trigger unsafe laboratory actions, or flood peer review with plausible but low-value manuscripts. In each case, the harm is not only a mistaken output. It is a distortion of the process by which a community decides what is true.

The current safety vocabulary is not wrong, but it is incomplete. Model cards, data sheets, red-team reports, and risk-management frameworks have improved transparency for AI systems \citep{mitchell2019modelcards,gebru2021datasheets,nist2023airmf}. Yet scientific AI introduces three features that require a more specific account. First, scientific systems operate on \emph{evidence loops}: their outputs may become data, citations, benchmarks, protocols, and future training material. Second, research agents are often connected to \emph{actuators}: code execution, cloud budgets, laboratory robots, procurement, synthesis providers, or publication channels. Third, science is a \emph{social epistemic system}: credibility depends on norms of attribution, replication, review, dissent, and public-interest justification. A deployment that optimizes for papers, novelty, or speed may improve local productivity while weakening the global reliability of science.

We therefore propose \position{} as a design and evaluation principle. Acceleration is desirable when it increases the rate at which communities produce reliable, useful, and responsibly governed knowledge. Acceleration is not desirable when it merely increases the rate of plausible artifacts, unreviewable claims, unsafe experiments, or unequal concentration of scientific power. Our concrete proposal is a \safetycase{}: a structured, evidence-backed argument that a scientific agent is acceptably safe and epistemically reliable for a specified deployment context. Safety cases are common in safety-critical engineering and are increasingly proposed for AI systems \citep{habli2025big}. We adapt the idea to the distinctive risks of scientific discovery.

\paragraph{Contributions.} This position paper makes four contributions. First, we introduce a ladder of scientific agency that clarifies how risk scales from drafting assistance to institutional discovery systems. Second, we identify five deployment risks for scientific AI: epistemic pollution, tool-mediated harm, responsibility laundering, inequitable concentration, and incentive distortion. Third, we propose a \safetycase{} framework with auditable claims about capability boundaries, tool containment, provenance, human accountability, and societal impact. Fourth, we outline a research agenda for LM4Sci: benchmarks and governance artifacts that measure not only whether agents can do science, but whether their science remains trustworthy.

\section{The Scientific Agency Gap}

A scientific language model is not automatically an autonomous scientist. However, scientific AI systems increasingly occupy roles that resemble parts of the scientific method: proposing hypotheses, designing interventions, interpreting results, and communicating claims. The resulting governance challenge is a gap between \emph{where agency appears} and \emph{where responsibility remains assigned}. Human researchers remain responsible for scientific claims, yet they may increasingly rely on agents whose search paths, tool calls, and evidential shortcuts are opaque, long-horizon, or difficult to audit.

Table~\ref{tab:ladder} presents a ladder of scientific agency. The levels are not meant as a moral ranking. They distinguish deployment contexts that require different evidence and safeguards. A read-only literature assistant may primarily require citation verification and privacy controls. A chemistry agent with access to synthesis planning tools requires dual-use screening, tool permissions, and expert oversight. An end-to-end research agent that produces papers requires provenance, reviewer protection, replication support, and policies against artificial inflation of publication volume.

\begin{table*}[t]
\centering
\small
\caption{A ladder of scientific agency. Risk increases not only with model capability but with autonomy, tool access, and institutional impact.}
\label{tab:ladder}
\begin{tabularx}{\textwidth}{@{}p{0.10\textwidth}p{0.24\textwidth}p{0.30\textwidth}X@{}}
\toprule
Level & Role & Typical action & Governance condition \\
\midrule
0. Text aid & Scientific autocomplete & Edits prose, code comments, or equations. & Authorship disclosure norms; no unverified claims. \\
1. Epistemic assistant & Retrieval and explanation & Finds papers, summarizes evidence, drafts related work. & Citation verification; uncertainty reporting; data privacy. \\
2. Workflow copilot & Bounded computation & Writes code, cleans data, runs notebooks, plots results. & Reproducible environments; sandboxing; human review of outputs. \\
3. Experimental operator & Bounded tool use & Executes simulations, queries databases, plans protocols under constraints. & Least-privilege permissions; budget limits; logged tool calls. \\
4. Research agent & Closed-loop discovery & Generates hypotheses, chooses experiments, analyzes failures, writes reports. & Safety case; independent replication; audit trail; escalation gates. \\
5. Institutional actor & Research program shaper & Allocates resources, submits papers, influences agendas or policy. & Public-interest governance; accountability board; incident reporting. \\
\bottomrule
\end{tabularx}
\end{table*}

The gap is easiest to miss because each step can look incremental. A literature assistant becomes a planner by adding search and memory. A planner becomes an experimental operator by adding a code interpreter. An experimental operator becomes a research agent by adding iterative hypothesis search. A research agent becomes an institutional actor when its outputs affect grant proposals, hiring, peer review, intellectual property, regulatory submissions, or public health. At no single step must the model become generally intelligent for the surrounding institution to change. The system becomes influential because it is embedded in workflows that already convert claims into resources and action.

This framing also clarifies why ``human in the loop'' is insufficient as a default safety answer. A human may be nominally present while lacking time, domain expertise, incentives, or trace access to evaluate the agent's work. In science, the relevant question is not whether a human clicked approve. It is whether a responsible human or institution understood enough about the agent's evidence, uncertainty, and hazards to stand behind the resulting claim. We therefore replace generic human-in-the-loop language with \emph{decision gates}: explicit points at which accountable humans authorize problem framing, tool escalation, resource expenditure, hazardous-domain work, claim release, and public communication.

\section{Risks Specific to Scientific AI Deployment}

The risks below overlap with general-purpose AI risks, but their scientific expression is distinctive. Science is cumulative. A small error can become a citation; a citation can become a benchmark; a benchmark can become a target; a target can reshape a field. The problem is not merely that agents may hallucinate. It is that scientific communities may inherit machine-generated artifacts without being able to tell which parts are reliable.

\subsection{Epistemic pollution of the scientific commons}

Scientific AI systems can reduce friction in producing plausible research artifacts. This is valuable when artifacts are correct, reproducible, and useful. It is dangerous when they are noisy, unverifiable, or optimized for acceptance rather than truth. The risk has several forms: fabricated citations, subtly corrupted metadata, benchmark contamination, overconfident summaries, non-reproducible code, synthetic datasets that encode hidden assumptions, and papers whose surface form resembles scientific contribution without adding durable knowledge. Recent large-scale audits of citations suggest that hallucinated or non-existent references are already a measurable problem in scientific writing workflows \citep{zhao2026hallucinations}. Autonomous paper-writing systems intensify this concern because they can generate not only sentences but entire chains of claims, plots, experiments, and reviews.

The harm is asymmetric. A false claim can be produced cheaply but may require expert time to detect and correct. If agentic systems increase submission volume faster than reviewing capacity, the scientific commons faces an epistemic denial-of-service problem: genuine reviewers spend more time filtering plausible artifacts and less time engaging deep contributions. This risk is especially salient for workshops and preprint servers, where speed and openness are strengths but also attack surfaces for low-quality automated production.

\subsection{Tool-mediated and dual-use harm}

Scientific agents differ from general chatbots because they often connect to tools that change the world. In computational settings, an agent may launch expensive jobs, scrape restricted data, or modify code repositories. In biological, chemical, or materials settings, an agent may design molecules, optimize protocols, suggest synthesis routes, or interact with robotic labs. Work on chemistry agents demonstrates the promise of tool-augmented systems \citep{bran2023chemcrow,boiko2023emergent}; the same architecture raises the stakes of misuse, accident, and inappropriate generalization.

Dual-use concerns are not limited to malicious users. A well-intentioned agent can create harm through specification errors, out-of-distribution reasoning, missing tacit knowledge, or over-optimization of a proxy. For example, a biological agent optimized to maximize novelty might search toward hazardous capabilities unless constrained by expert screening. A materials agent might suggest experiments that are chemically unsafe under real laboratory conditions. A code agent might generate insecure pipelines for handling sensitive biomedical data. The International Scientific Report on advanced AI safety emphasizes misuse, malfunction, controllability, and systemic risks for general-purpose AI \citep{bengio2024international}; scientific agents make these concerns concrete by tying models to domain-specific tools.

\subsection{Responsibility laundering}

Scientific responsibility currently attaches to people and institutions: authors, principal investigators, laboratories, companies, universities, journals, funders, and regulators. Autonomous agents complicate this structure. When an agent proposes a hypothesis, chooses an analysis, and drafts the paper, who is responsible for the evidential standard? The developer who trained the model? The scientist who deployed it? The lab that provided tools? The reviewer who accepted the claim? The answer cannot be ``the AI,'' because the system cannot assume moral, legal, or professional accountability.

Responsibility laundering occurs when AI involvement diffuses blame while preserving credit. A lab may credit the system for speed and novelty when outcomes are positive, but treat failures as unpredictable model behavior when harms occur. This is especially concerning in science because credit allocation influences careers, funding, and trust. The norm should be simple: AI systems may be acknowledged, but accountable humans and institutions must own the claims and the deployment conditions that produced them.

\subsection{Inequality and concentration}

Scientific AI could democratize expertise by giving smaller labs access to tools for coding, literature analysis, and experimental design. It could also concentrate scientific power. Frontier models, proprietary training data, private laboratory automation, and large compute budgets may allow a small number of organizations to shape research agendas. If the most capable scientific agents are closed systems, the broader community may be asked to trust outputs it cannot inspect or reproduce. AlphaFold demonstrated the immense public value of releasing predictions and tools that reshape a field \citep{jumper2021alphafold}; debates around later systems also show that access conditions, code release, and commercial incentives matter for scientific legitimacy \citep{abramson2024alphafold3}.

The distributional question is not only who can use agents. It is whose problems agents are optimized to solve. Systems trained on the published literature may overweight well-funded topics, English-language papers, positive results, and historically dominant institutions. If deployed at scale, agents could amplify existing inequities in recognition, agenda setting, and resource allocation. A responsible LM4Sci agenda must therefore evaluate benefits and burdens across researchers, disciplines, countries, and affected communities.

\subsection{Incentive distortion and loss of scientific judgment}

Scientific agents will be deployed inside incentive systems that reward novelty, publication volume, citation impact, benchmark performance, patents, and speed. Without safeguards, agents may optimize these proxies. An agent that can cheaply generate many plausible experiments may encourage shallow search rather than hard conceptual work. A paper-writing agent may make weak results look polished. A literature agent may converge on fashionable framings. A reviewer agent may reward conventional argument structure over genuine insight. The result is not simply automation bias; it is a shift in what the institution learns to value.

This risk has a human dimension. Scientific judgment is a skill developed through failed experiments, tacit laboratory knowledge, peer disagreement, and conceptual struggle. Offloading every difficult step to agents may deskill researchers or narrow their sense of what counts as evidence. Conversely, well-designed agents can improve judgment by exposing uncertainty, suggesting alternative hypotheses, and making hidden assumptions legible. The design goal should be augmentation of scientific understanding, not replacement of the reflective practices that make science self-correcting.

\section{Position: Accountable Acceleration}

The promise of LM4Sci is real. Language models can lower barriers to entry, connect distant literatures, automate tedious analysis, and search spaces too large for unaided humans. The position of this paper is not to slow scientific AI by default. It is to reject a one-dimensional notion of speed. We should ask: speed toward what, under whose control, with what evidence, and at whose risk?

We define \position{} as the principle that scientific AI should increase the rate of reliable discovery while preserving the social conditions that make discovery trustworthy. Four commitments follow.

\paragraph{Acceleration should be measured at the level of validated knowledge, not artifact production.} A system that produces ten papers with weak evidence has not accelerated science in the same way as a system that produces one robust, replicable result. LM4Sci evaluations should therefore include replication, claim verification, negative-result handling, and downstream expert usefulness, not only novelty or paper acceptance.

\paragraph{Scientific agents should be permissioned by action, not by model identity.} The same base model can be low risk in a read-only literature setting and high risk when connected to synthesis tools or publication pipelines. Governance should track capabilities, tools, domains, and user roles. A chemistry agent with procurement access should face stronger controls than a summarization assistant, even if both use the same language model.

\paragraph{Provenance should be a first-class research artifact.} A scientific claim produced with AI assistance should be traceable to data, code, prompts, tool calls, model versions, intermediate failures, human interventions, and validation checks. Not every private detail can be public, especially in sensitive or proprietary research. But the default should shift from polished final paper to evidence graph: what was tried, what failed, what changed, and who approved each step.

\paragraph{Human accountability should be located at decision gates.} Oversight must be tied to consequential actions: choosing the research question, escalating tool permissions, spending resources, entering dual-use domains, making claims public, and deploying outputs. The accountable human is not a passive monitor. They are the person or institution that can justify why the agent's operation was appropriate, what evidence supports its outputs, and how residual risks are managed.

\section{A Scientific-Agent Safety Case}

A safety case is a structured argument, supported by evidence, that a system is acceptably safe for a defined context. For scientific agents, the argument must cover both conventional safety and epistemic integrity. Table~\ref{tab:safetycase} sketches the proposed structure. The point is not to create paperwork for its own sake. The point is to make deployment claims inspectable before harm occurs and revisable as systems learn, tools change, or failures appear.

\begin{table*}[t]
\centering
\small
\caption{A proposed \safetycase{} structure for autonomous research agents. Each claim should be backed by evidence, assumptions, and known defeaters.}
\label{tab:safetycase}
\begin{tabularx}{\textwidth}{@{}p{0.19\textwidth}p{0.37\textwidth}X@{}}
\toprule
Safety-case claim & Required evidence & Example defeaters \\
\midrule
Purpose and scope & Intended domain, users, excluded uses, risk tier, benefit claim. & Vague objective; hidden commercial or policy use; unbounded autonomy. \\
Capability boundaries & Domain evals, calibration tests, red-team results, failure-mode analysis. & Benchmark contamination; brittle long-horizon planning; overconfident uncertainty. \\
Tool containment & Least-privilege permissions, sandboxing, budget limits, approval gates, emergency stop. & Unlogged tool calls; indirect access to prohibited tools; prompt-injection path. \\
Epistemic provenance & Data/code versions, claim-to-evidence graph, citation checks, negative results, replication package. & Missing intermediate failures; unverifiable citations; irreproducible environment. \\
Human accountability & Named accountable role, review protocol, escalation criteria, incident response. & Rubber-stamp approval; unclear ownership across developer/deployer/lab. \\
Societal impact & Access analysis, affected stakeholders, dual-use review, environmental/resource costs. & Benefits captured by few actors; unsafe open release; unexamined downstream misuse. \\
\bottomrule
\end{tabularx}
\end{table*}

\subsection{Capability boundaries}

A scientific agent should have a stated operating envelope. This includes domains, tasks, assumptions, model versions, known failure modes, and conditions under which the agent must refuse or escalate. Capability evaluations should be domain-specific and adversarial. For literature agents, evaluation should include citation existence, claim support, disagreement retrieval, and handling of retracted or low-quality sources. For data-analysis agents, evaluation should include leakage detection, statistical validity, robustness to mislabeled data, and reproducibility. For experimental agents, evaluation should include protocol safety, uncertainty propagation, and performance under failed or anomalous experiments.

Crucially, the boundary is not a single benchmark score. It is an argument about why the system is reliable enough for a deployment. Benchmarks can be part of the evidence, but they should be accompanied by stress tests, expert review, and explicit residual uncertainty. This is consistent with the broader movement from abstract trustworthiness principles to context-sensitive AI risk management \citep{nist2023airmf,nist2024genai}.

\subsection{Tool containment and graduated permissions}

Scientific agents should follow least privilege. A literature agent should not inherit the permissions of a laboratory robot. A code agent should not automatically obtain network access, cloud spending authority, or the ability to submit manuscripts. We recommend a permission stack: read, write, execute, spend, actuate, communicate, and publish. Each level should require separate authorization, monitoring, and revocation.

Tool containment should include sandboxed execution, package allowlists, data egress controls, budget caps, rate limits, prompt-injection defenses, and human approval for hazardous domains. For wet-lab or dual-use settings, the agent should not directly order materials, design restricted biological constructs, or execute dangerous protocols without institutional biosafety or chemical-safety review. A practical rule is that the agent should never be the only entity that understands why a potentially hazardous action is acceptable.

\subsection{Epistemic provenance}

Scientific AI should make the path from hypothesis to claim inspectable. We propose a provenance triad: \emph{data provenance} records datasets, instruments, preprocessing, and licenses; \emph{reasoning provenance} records prompts, plans, tool outputs, uncertainty estimates, and rejected alternatives; \emph{intervention provenance} records human edits, approvals, escalations, and overrides. The goal is not to preserve every token forever. The goal is to preserve enough structure for reviewers, collaborators, auditors, and future scientists to understand how the claim was produced.

Provenance also changes incentives. If agents must record failed hypotheses and negative results, they become less useful for p-hacking and more useful for cumulative science. If papers include claim-to-evidence links, reviewers can inspect the weakest parts of an argument rather than infer them from polished prose. If generated citations are machine-checked against scholarly databases, fabricated references become less likely to enter the record. These mechanisms are technical, but their purpose is institutional: protect the commons of evidence.

\subsection{Accountability and incident response}

Every deployment of a nontrivial scientific agent should identify accountable roles. At minimum, there should be a deploying scientist responsible for claims, a technical owner responsible for system behavior, and an institutional owner responsible for risk acceptance. For high-risk domains, an independent review function should approve the safety case before tool escalation.

Incident response should be explicit. What happens if an agent fabricates a citation, leaks data, generates a hazardous protocol, publishes an erroneous result, or corrupts a benchmark? Who is notified? Are outputs retracted? Are downstream users warned? Are logs preserved? Does the agent lose permissions? Scientific AI governance should borrow from cybersecurity and laboratory safety: incidents are not shameful anomalies but expected signals for improving controls.

\section{Implications for Evaluation and Publication}

The LM4Sci community can shape norms before autonomous research agents become ordinary infrastructure. We propose four changes to evaluation and publication practice.

\paragraph{Evaluate workflows, not only answers.} A system that reaches a correct answer through untrustworthy steps may be unsafe to deploy. Evaluation should include process traces, tool-call audits, failed-run analysis, and human oversight burden. For research agents, a benchmark should ask not only whether the final paper looks plausible, but whether the experiments were meaningful, the claims supported, the code reproducible, and the negative results preserved.

\paragraph{Require AI-use and provenance statements for agentic submissions.} COLM submission instructions already require disclosure when LLMs originate ideas, write original content, generate data or plots, or perform evaluation \citep{colm2026instructions}. LM4Sci should go further for agentic work: submissions should disclose what parts of the research loop were automated, what tools the system could access, what human gates existed, and how citations and empirical claims were verified.

\paragraph{Protect peer review from automated volume.} Workshops should welcome work on automated discovery while resisting incentives that reward paper factories. Possible safeguards include mandatory citation validation, reproducibility checklists, limits on fully automated submissions without human-authored critical analysis, reviewer-facing provenance summaries, and explicit rejection criteria for unverified generated content. The goal is not to ban AI-authored drafts; it is to prevent review capacity from becoming the bottleneck that low-quality automation exploits.

\paragraph{Reward negative results and safety failures.} Scientific-agent papers should be encouraged to report unsafe behaviors, failed hypotheses, misleading experiments, and oversight breakdowns. These are not embarrassing side notes; they are core evidence for responsible deployment. A workshop on language models for science is an ideal venue for normalizing failure reports before the field becomes too performance-centric.

\section{A Research Agenda for LM4Sci}

We close with open problems that are technical enough for research papers and institutional enough to matter.

\paragraph{Agentic scientific safety benchmarks.} We need benchmarks that evaluate long-horizon scientific workflows under realistic constraints: incomplete literature, noisy data, failed experiments, ambiguous hypotheses, and tool hazards. Benchmarks should include safe-refusal and escalation behavior, not only task success. They should measure whether agents notice uncertainty, seek disconfirming evidence, and avoid unsafe tool use.

\paragraph{Claim verification and evidence graphs.} Scientific text should be linked to machine-checkable evidence where possible. A claim graph connects each contribution to supporting data, code, citations, and human decisions. Research is needed on usable formats, privacy-preserving logs, selective disclosure, cryptographic signing, and reviewer interfaces that reduce rather than increase burden.

\paragraph{Oversight economics.} Human oversight is not free. A nominally safe system that requires experts to inspect thousands of tool calls may be unusable. We need empirical studies of oversight load: how long experts need to audit agent traces, which summaries preserve understanding, when automation bias appears, and which decision gates actually catch failures.

\paragraph{Dual-use evaluation with domain experts.} Scientific AI safety cannot be solved by generic LLM red teaming alone. Biology, chemistry, materials science, cybersecurity, and medicine require domain-specific threat models. The community should develop controlled evaluation protocols that involve experts without broadly disseminating dangerous details.

\paragraph{Governance of the scientific commons.} Scientific AI depends on public literature, open-source code, shared databases, and peer review. The community should study reciprocal obligations: when systems train on or retrieve from the commons, what should they return? Possible answers include verified metadata, replication artifacts, negative-result repositories, open evaluation logs, and funding for shared infrastructure.

\paragraph{Norms for AI participation in authorship and credit.} AI systems should not be listed as accountable authors, but their contributions should not be hidden. We need norms that distinguish editing, analysis, hypothesis generation, experiment selection, and manuscript drafting. Such norms should protect early-career scientists from both unfair suspicion and unfair displacement.

\section{Conclusion}

The most important question for scientific AI is not whether autonomous agents can produce papers. They can, and their capabilities will improve. The question is whether they can participate in science without eroding the conditions that make science trustworthy. A field that optimizes only for speed may discover more artifacts while knowing less. A field that treats scientific agents as accountable infrastructure can instead pursue \position{}: faster discovery with stronger provenance, clearer responsibility, safer tool use, and fairer access.

LM4Sci is well positioned to define this agenda. Its community sits at the intersection of language modeling, scientific practice, evaluation, and human-AI collaboration. The workshop should therefore ask every proposed system a simple set of questions: What decisions does the agent make? What tools can it use? What evidence supports its claims? Who is accountable? What happens when it fails? These questions should not be seen as obstacles to discovery. They are the infrastructure that will allow discovery to accelerate without losing public trust.

\section*{Ethics Statement}
This paper argues for stronger governance of scientific AI systems and does not introduce a new model, dataset, or experimental capability. The main ethical risk of the paper is misuse of the framing as compliance theater: organizations might produce superficial safety cases without changing deployment behavior. To reduce that risk, we emphasize evidence-backed claims, known defeaters, independent review for high-risk domains, and incident response. We also note that safety controls can be misused to centralize power or restrict beneficial open science; therefore governance should include access, equity, and public-interest considerations.

\section*{Disclosure of LLM Assistance}
A language model was used to assist with drafting and editing this position-paper text. The submitting authors are responsible for verifying all claims, references, and policy statements before submission. This disclosure is included because COLM guidance requires disclosure when LLMs are used to write original paper content.

\bibliography{lm4sci_references}
\bibliographystyle{colm2026_conference}

\appendix
\section{Reviewer-Facing Checklist for a Scientific-Agent Safety Case}
\label{app:checklist}
\small
\begin{enumerate}
  \item \textbf{Scope.} Does the paper specify the agent's intended users, domain, autonomy level, and excluded uses?
  \item \textbf{Tools.} Does it list the tools the agent can access, with permissions and containment mechanisms?
  \item \textbf{Evidence.} Are major scientific claims linked to data, code, citations, or experiments that reviewers can inspect?
  \item \textbf{Verification.} Were citations, generated data, plots, and code independently checked?
  \item \textbf{Oversight.} Which decisions required human approval, and what expertise did the approving humans have?
  \item \textbf{Failure reporting.} Does the paper report failed runs, unsafe suggestions, hallucinations, or other negative evidence?
  \item \textbf{Risk tier.} Does the deployment involve biological, chemical, medical, cyber, environmental, or other high-impact domains?
  \item \textbf{Societal impact.} Who benefits, who bears risk, and who can reproduce or contest the output?
\end{enumerate}

\end{document}
